\documentclass[11pt]{article}
\usepackage{fontspec,amsmath,amssymb,amsthm,hyperref}
\usepackage[margin=1in]{geometry}
\setlength{\emergencystretch}{3em}
\title{An independent Brownian family for all actual particle coordinates}
\author{LangevinLean mathematical exposition}
\date{}
\begin{document}\maketitle
\noindent Source module: \texttt{ProximalLangevin/MeanFieldBrownianFamily.lean}. Compile this proof document with XeLaTeX or LuaLaTeX. The accompanying HTML pairs the mathematical argument with the exact Lean statements and records the scope of the formalization.


On \(Q^{\otimes N}\), choose parameter Brownian coordinates from copy zero and latent coordinate \((i,j)\) from copy \(i\). Injective selection proves independence of complete selected coordinate paths.


\section{Select distinct copy/coordinate pairs}


Map a parameter index \(j\) to \((0,\mathrm{parameter}\ j)\), and a latent index \((i,j)\) to \((i,\mathrm{latent}\ j)\). The parameter/latent tags make the two kinds disjoint, and equality within either kind forces equality of the original indices. This proves injectivity. Define the family by evaluating the corresponding source Brownian coordinate in the selected copy; the displayed coordinate identities unfold the selection.




The index map sends a parameter coordinate \(j\) to \((0,\theta_j)\) and a latent particle coordinate \((i,j)\) to \((i,x_j)\). These are distinct even for the first particle because the parameter and latent coordinates use different summands of the joint index. The injectivity proof checks the four parameter/latent cases explicitly. The two evaluation lemmas then identify the selected paths with the intended common parameter noise and each particle's own latent noise.




\textbf{Correspondence with the displayed declarations.}


\begin{itemize}

\item \texttt{particleBrownianIndex}: Select a distinct base-copy/coordinate pair for every particle-noise coordinate.
\item \texttt{particleBrownianIndex\_injective}: The common parameter uses copy zero; latent blocks use their respective copies.
\item \texttt{particleBrownianFamily}: The actual joint Brownian source for the interacting particle coordinates.
\end{itemize}


\noindent Lean declarations: \texttt{particleBrownianIndex}, \texttt{particleBrownianIndex\_injective}, \texttt{particleBrownianFamily}, \texttt{particleBrownianFamily\_parameter}, \texttt{particleBrownianFamily\_latent}.

\section{Each selected process retains its Brownian law}


Pullback along product evaluation is measure preserving, so every selected scalar process is standard pre-Brownian. If all source evaluations are measurable, composition with the selected coordinate map proves every new evaluation measurable as well.




\textbf{Correspondence with the displayed declarations.}


\begin{itemize}

\item \texttt{particleBrownianFamily\_preBrownian}: Every selected coordinate is an actual standard pre-Brownian motion.
\item \texttt{particleBrownianFamily\_measurable}: Measurable source evaluations remain measurable after selecting a copy.
\end{itemize}


\noindent Lean declarations: \texttt{particleBrownianFamily\_preBrownian}, \texttt{particleBrownianFamily\_measurable}.

\section{Flatten independent rows and select an injective subfamily}


The finite product makes the complete Brownian rows independent. Inside each row, the supplied full coordinate paths are independent. The product-family independence theorem flattens these into independent variables indexed by \((i,j)\). Precomposing with the injective selection proves independence of all particle coordinates, including parameter and first-particle latent blocks. The canonical Brownian source supplies both the row laws and full-path independence, yielding an explicit complete family.




First pull each whole Brownian row from the source to one coordinate of the finite product; distinct rows are independent by the product measure. Within a row the supplied coordinate paths are independent, and measure-preserving evaluation retains their joint law. The helper \texttt{iIndepFun\_uncurry'} flattens these two levels into a jointly independent family indexed by all copy/coordinate pairs. Finally, precomposition with the injective index map selects the intended subfamily and preserves independence. This proves independence of full coordinate paths, including the parameter path and the first particle's latent path; it does not assert independence of interacting particle states.




\textbf{Correspondence with the displayed declarations.}


\begin{itemize}

\item \texttt{particleBrownianFamily\_paths\_independent}: Full coordinate-path independence is derived on the actual finite product. There is no assumption of independent particles after interaction.
\item \texttt{canonical\_particleBrownianFamily}: The concrete canonical Gaussian time-indicator construction supplies all selected scalar Brownian laws and their full-path independence.
\end{itemize}


\noindent Lean declarations: \texttt{particleBrownianFamily\_paths\_independent}, \texttt{canonical\_particleBrownianFamily}.

\end{document}
